\documentclass[10pt,a4paper]{amsart}
\usepackage{amsmath,amssymb,amsthm,graphicx}
\usepackage{verbatim}
\newenvironment{reference}{\comment}{\endcomment}
%\newenvironment{reference}{}{}
\newenvironment{slogan}{\comment}{\endcomment}
\newenvironment{history}{\comment}{\endcomment}
% For commutative diagrams we use Xy-pic
\usepackage[all]{xy}

% We use 2cell for 2-commutative diagrams.
\xyoption{2cell}
% Theorem environments.
%
\theoremstyle{plain}
\newtheorem{theorem}[subsection]{Theorem}
\newtheorem{proposition}[subsection]{Proposition}
\newtheorem{lemma}[subsection]{Lemma}

\theoremstyle{definition}
\newtheorem{definition}[subsection]{Definition}
\newtheorem{example}[subsection]{Example}
\newtheorem{exercise}[subsection]{Exercise}
\newtheorem{situation}[subsection]{Situation}

\theoremstyle{remark}
\newtheorem{remark}[subsection]{Remark}
\newtheorem{remarks}[subsection]{Remarks}

\numberwithin{equation}{subsection}

% Macros
%
\def\lim{\mathop{\mathrm{lim}}\nolimits}
\def\colim{\mathop{\mathrm{colim}}\nolimits}
\def\Spec{\mathop{\mathrm{Spec}}}
\def\Hom{\mathop{\mathrm{Hom}}\nolimits}
\def\Ext{\mathop{\mathrm{Ext}}\nolimits}
\def\SheafHom{\mathop{\mathcal{H}\!\mathit{om}}\nolimits}
\def\SheafExt{\mathop{\mathcal{E}\!\mathit{xt}}\nolimits}
\def\Sch{\mathit{Sch}}
\def\Mor{\mathop{\mathrm{Mor}}\nolimits}
\def\Ob{\mathop{\mathrm{Ob}}\nolimits}
\def\Sh{\mathop{\mathit{Sh}}\nolimits}
\def\NL{\mathop{N\!L}\nolimits}
\def\CH{\mathop{\mathrm{CH}}\nolimits}
\def\proetale{{pro\text{-}\acute{e}tale}}
\def\etale{{\acute{e}tale}}
\def\QCoh{\mathit{QCoh}}
\def\Ker{\mathop{\mathrm{Ker}}}
\def\Im{\mathop{\mathrm{Im}}}
\def\Coker{\mathop{\mathrm{Coker}}}
\def\Coim{\mathop{\mathrm{Coim}}}

% Boxtimes
%
\DeclareMathSymbol{\boxtimes}{\mathbin}{AMSa}{"02}

%
% Macros for moduli stacks/spaces
%
\def\QCohstack{\mathcal{QC}\!\mathit{oh}}
\def\Cohstack{\mathcal{C}\!\mathit{oh}}
\def\Spacesstack{\mathcal{S}\!\mathit{paces}}
\def\Quotfunctor{\mathrm{Quot}}
\def\Hilbfunctor{\mathrm{Hilb}}
\def\Curvesstack{\mathcal{C}\!\mathit{urves}}
\def\Polarizedstack{\mathcal{P}\!\mathit{olarized}}
\def\Complexesstack{\mathcal{C}\!\mathit{omplexes}}
% \Pic is the operator that assigns to X its picard group, usage \Pic(X)
% \Picardstack_{X/B} denotes the Picard stack of X over B
% \Picardfunctor_{X/B} denotes the Picard functor of X over B
\def\Pic{\mathop{\mathrm{Pic}}\nolimits}
\def\Picardstack{\mathcal{P}\!\mathit{ic}}
\def\Picardfunctor{\mathrm{Pic}}
\def\Deformationcategory{\mathcal{D}\!\mathit{ef}}
\usepackage{cleveref}
\crefname{originalreference}{original reference}{original references}
\Crefname{originalreference}{Original reference}{Original references}
\makeatletter
\newlabel{dualizing-lemma-local-cohomology-adjoint}{{dualizing-lemma-local-cohomology-adjoint}{source}{}{sourcepointer-1}{}}
\newlabel{dualizing-lemma-local-cohomology-adjoint@cref}{{[originalreference][0][]dualizing-lemma-local-cohomology-adjoint}{[1][0][]source}{}{}{}}
\makeatother
\crefname{originalreference}{original reference}{original references}
\Crefname{originalreference}{Original reference}{Original references}
\makeatletter
\newlabel{more-algebra-remark-P-resolution}{{more-algebra-remark-P-resolution}{source}{}{sourcepointer-2}{}}
\newlabel{more-algebra-remark-P-resolution@cref}{{[originalreference][0][]more-algebra-remark-P-resolution}{[1][0][]source}{}{}{}}
\makeatother
\title{The derived sections with support are computed by the derived tensor product with the sections of the structure sheaf, in finite dimension}
\author{The Stacks Project Authors}
\date{stacks-project:local-cohomology.tex:local-snapshot-20261012; GNU FDL 1.2 or later}
\begin{document}
\maketitle
\noindent\textit{Editorial note.} All delivered mathematical text is the author's own, copied from the pinned Stacks Project chapter source sources/stacks-local-cohomology/pinned-original-source.tex; no statement or proof step was rewritten, reordered or removed.\par\medskip
\noindent\textit{Editorial note.} Preamble provenance: the Stacks Project keeps one single shared preamble for every chapter; this card takes it directly from the staged sources/stacks-local-cohomology/preamble.tex as three verbatim spans (lines 8-12, 14-18 and 154-217), so no derived file was written and nothing was added or redefined. The wrapper class and the base packages are supplied by the assembler.\par\medskip
\noindent\textit{Editorial note.} Bibliography provenance: neither the counted statement nor its proof contains a citation (this lane rescans the whole span before writing the card), so the card carries no bibliography and no entry was reconstructed.\par\medskip
\noindent\textit{Editorial note.} Selection provenance: this card was opened under the file-based control contract (CONTROL.md), with the mechanical pre-card gate run on its statement and proof interval before the card was written.\par\medskip
\noindent\textit{Allowed standard prerequisites (editorial).}\begin{itemize}
\item The description of the sections with support in a specialization-stable subset as the filtered colimit of the sections with support in the closed subsets it contains. Author lemma lemma-adjoint-ext, whose statement is reproduced verbatim as a context span of this card (source lines 996-1004). Combined with the commutativity of directed colimits with direct sums, gives that the functor commutes with direct sums so that the map may be checked on shifts of the ring.
\item K-flat complexes as a model for the derived tensor product and the total complex of a double complex. Used at the start of the counted proof (source lines 1180-1184). Provides the explicit representative of the derived tensor product on which the canonical map is built.
\item The reduction of an isomorphism of functors commuting with direct sums to the shifts of the ring. Invoked inside the counted proof at source lines 1196-1198 as a cross-chapter pointer. Reduces the verification of the displayed identity to the case of a shift of the ring, where it is immediate.
\end{itemize}\medskip
\ifdefined\hypertarget\hypertarget{sourcepointer-1}{}\fi
\noindent\textit{Source reference.} Original reference dualizing-lemma-local-cohomology-adjoint: sources/stacks-local-cohomology/pinned-original-source.tex, source lines 335--335. Navigation pointer only: the counted text refers to the Stacks item dualizing-lemma-local-cohomology-adjoint, which belongs to another chapter of the project and is not part of this card. The pointer carries the target's own label as its number because this chapter's staged snapshot does not contain the other chapter's tag table; its content is not reproduced here.\par\medskip
\ifdefined\hypertarget\hypertarget{sourcepointer-2}{}\fi
\noindent\textit{Source reference.} Original reference more-algebra-remark-P-resolution: sources/stacks-local-cohomology/pinned-original-source.tex, source lines 1196--1196. Navigation pointer only: the counted text refers to the Stacks item more-algebra-remark-P-resolution, which belongs to another chapter of the project and is not part of this card. The pointer carries the target's own label as its number because this chapter's staged snapshot does not contain the other chapter's tag table; its content is not reproduced here.\par\medskip
\begin{lemma}
\label{lemma-torsion-tensor-product}
Let $A$ be a Noetherian ring. Let $T \subset \Spec(A)$ be a subset
stable under specialization. Assume $A$ has finite dimension. Then
$$
R\Gamma_T(K) = R\Gamma_T(A) \otimes_A^\mathbf{L} K
$$
for $K \in D(A)$. For $K, L \in D(A)$ we have
$$
R\Gamma_T(K \otimes_A^\mathbf{L} L) =
K \otimes_A^\mathbf{L} R\Gamma_T(L) =
R\Gamma_T(K) \otimes_A^\mathbf{L} L =
R\Gamma_T(K) \otimes_A^\mathbf{L} R\Gamma_T(L)
$$
If $K$ or $L$ is in $D_T(A)$ then so is $K \otimes_A^\mathbf{L} L$.
\end{lemma}

\begin{lemma}
\label{lemma-adjoint-ext}
Let $A$ be a Noetherian ring. Let $T \subset \Spec(A)$
be a subset stable under specialization.
For any object $K$ of $D(A)$ we have
$$
H^i(RH^0_T(K)) = \colim_{Z \subset T\text{ closed}} H^i_Z(K)
$$
\end{lemma}

\begin{proof}
By construction we may represent $R\Gamma_T(A)$ by a complex $J^\bullet$ in
$\text{Mod}_{A, T}$. Thus if we represent $K$ by a K-flat complex $K^\bullet$
then we see that $R\Gamma_T(A) \otimes_A^\mathbf{L} K$ is represented
by the complex $\text{Tot}(J^\bullet \otimes_A K^\bullet)$ in
$\text{Mod}_{A, T}$. Using the map $R\Gamma_T(A) \to A$ we obtain
a map $R\Gamma_T(A) \otimes_A^\mathbf{L} K\to K$. Thus by the adjointness
property of $R\Gamma_T$ we obtain a canonical map
$$
R\Gamma_T(A) \otimes_A^\mathbf{L} K \longrightarrow R\Gamma_T(K)
$$
factoring the just constructed map. Observe that $R\Gamma_T$ commutes
with direct sums in $D(A)$ for example by Lemma \ref{lemma-adjoint-ext},
the fact that directed colimits commute with direct sums, and the
fact that usual local cohomology commutes with direct sums
(for example by Dualizing Complexes, Lemma
\ref{dualizing-lemma-local-cohomology-adjoint}).
Thus by More on Algebra, Remark \ref{more-algebra-remark-P-resolution}
it suffices to check the map is an isomorphism for
$K = A[k]$ where $k \in \mathbf{Z}$. This is clear.

\medskip\noindent
The final statements follow from the result we've just shown
by transitivity of derived tensor products.
\end{proof}

\end{document}
